% Complete independent proof fragment; requires amsmath, amssymb, hyperref.
% Distinct equation prefix HD. Original HH and HC sources are not modified.
\section{The actual heat equation annihilates the integral monodromy defect}
\label{sec:hd}
\newcounter{hd}\renewcommand{\thehd}{HD\arabic{hd}}
\newcommand{\HD}{\refstepcounter{hd}\par\medskip\noindent\textbf{\thehd.}\ }

\HD\textbf{The original objects and the exact question.}
Retain the original Rodgers--Tao convention
\[
 H_t(Z)=\int_0^\infty e^{tu^2}\Phi(u)\cos(Zu)\,du,\qquad
 H_0(Z)=\tfrac18\xi_{\mathrm R}(\tfrac12+iZ/2),\qquad
 \partial_tH_t=-\partial_Z^2H_t,
 \tag{HD1}\label{hd:heat}
\]
where
$\Phi(u)=\sum_{n\ge1}(2\pi^2n^4e^{9u}-3\pi n^2e^{5u})
e^{-\pi n^2e^{4u}}$ and
$\xi_{\mathrm R}(s)=s(s-1)\pi^{-s/2}\Gamma(s/2)\zeta(s)/2$.
The source is Brad Rodgers and Terence Tao,
\emph{The de Bruijn--Newman constant is non-negative},
\href{https://arxiv.org/abs/1801.05914v5}{arXiv:1801.05914v5},
original author equations \texttt{hoz}, \texttt{sas},
\texttt{phidef}, \texttt{htdef} and the following heat-equation sentence.
Fix an actual real zero $x_0$ of $H_0$, of its actual order $m\ge1$.
No multiple-zero existence is assumed. Write, in the original coordinate,
\[
 \xi=Z-x_0,\quad H_0(x_0+\xi)=\sum_{n\ge m}a_n\xi^n,\quad
 a_m\ne0,\qquad A_k=\frac{a_{m+k}}{a_m},\quad A_0=1.
 \tag{HD2}
\]
These real coefficients, including $a_m$, are retained throughout.
The double-exponential decay in HD1 permits arbitrary derivatives
on compact complex $(t,Z)$ sets, dominated by
$C(1+u)^N e^{C u^2+Cu-ce^{4u}}$ for $u\ge1$.
This proves joint holomorphy and legitimizes all local Taylor
identities below.

Put $\mathcal O=\mathbb C\{t\}$ and
$\mathcal K=\operatorname{Frac}\mathcal O$.
The exact monic polynomial $W(t,\xi)$ whose roots are this local
cluster gives
\[
 H_t(x_0+\xi)=U(t,\xi)W(t,\xi),\qquad
 A=\mathcal O[\xi]/W,\quad U(0,\xi)=\sum_{k\ge0}a_{m+k}\xi^k .
 \tag{HD3}\label{hd:original}
\]
The entire analytic unit $U$ is retained.
The auxiliary algebra $B$ and the inclusion $\iota:A\hookrightarrow B$
will be constructed explicitly below. On $B$, $S_B$ changes the
sign of each square-root coordinate. HH14 left the finite module
\[
 \mathcal D_M=(A+M_AA)/A\subset B/A,\qquad
 M_A=\iota_{\mathcal K}^{-1}S_B\iota_{\mathcal K}
 \tag{HD4}\label{hd:defect}
\]
without assuming its vanishing. The conclusion proved here is
$M_AA=A$ and $\mathcal D_M=0$ for every actual multiplicity $m$.
It follows from the heat identities for all the original
coefficients, rather than from a general property of finite
normalization or a presumption that holonomy preserves a lattice.

\HD\textbf{The full heat branches and their defining recursion.}
Define the monic polynomials
\[
 P_n(z)=\sum_{j=0}^{\lfloor n/2\rfloor}
      \frac{(-1)^jn!}{j!(n-2j)!}z^{n-2j},\qquad
 \mathcal P_m(t,y)=
      \sum_{j=0}^{\lfloor m/2\rfloor}
       \frac{(-1)^jm!}{j!(m-2j)!}t^jy^{m-2j}.
 \tag{HD5}\label{hd:hermite}
\]
This is exactly the convention
$P_n(z)=H_n^{\mathrm{Hermite}}(z/2)$ of
\href{https://dlmf.nist.gov/18.5.E13}{DLMF 18.5.13},
Chapter 18 by T.H. Koornwinder, R. Wong, R. Koekoek and
R.F. Swarttouw. The finite formulas fix all constants independently.
Their generating series is $e^{zw-w^2}$, giving
$P_n'=nP_{n-1}$,
$P_{n+1}=zP_n-2nP_{n-1}$, and
$2P_n''-zP_n'+nP_n=0$.
Also $P_ne^{-z^2/4}=(-2)^n\partial_z^n e^{-z^2/4}$, so integration
by parts makes $P_n$ orthogonal to each polynomial of smaller
degree for the positive Gaussian weight. If it had fewer than
$n$ real sign-changing roots, multiplication by the product of
their distinct linear factors would give a nonzero function of
constant sign and contradict that orthogonality. Thus its $n$
roots are real and simple.

Let $r_1<\cdots<r_m$ be the roots of $P_m$ and
$\sigma j=m+1-j$. Then $r_{\sigma j}=-r_j$.
The original heat Taylor identity is the convergent equality
\[
 H_{u^2}(x_0+uz)
       =a_mu^m\sum_{k\ge0}A_k u^kP_{m+k}(z),\qquad t=u^2.
 \tag{HD6}\label{hd:fullseries}
\]
Indeed the heat equation gives
$[t^j\xi^\ell]H_t=(-1)^j(\ell+2j)!a_{\ell+2j}/(j!\ell!)$.
Regroup the absolutely convergent Taylor series by $n=2j+\ell$
on a sufficiently small polydisc.
The analytic implicit-function theorem at $(u,z)=(0,r_j)$ gives
all the branches
\[
 \theta_j(u)=u z_j(u),\quad z_j(0)=r_j,\qquad
 \theta_j(-u)=\theta_{\sigma j}(u),\qquad
 \theta_j(\bar u)=\overline{\theta_j(u)} .
 \tag{HD7}
\]
An isolating circle and the argument principle show that these
are precisely the $m$ zeros of the original nearby cluster.
Every difference is $u(r_j-r_k+O(u))$, so they are distinct off
$u=0$. Their product gives $W$ in HD3.
The quotient $H/W$ is jointly analytic by successive holomorphic
division by the branches on the $u$ cover and invariance under
$u\mapsto-u$; its value at $(0,0)$ is $a_m$, proving that it is a unit.

\HD\textbf{A degree bound forced by the heat equation.}
For all integers $k,q\ge0$ there are polynomials $Q_{k,q}(z)$
such that at every root $r$ of $P_m$,
\[
 \frac{P_{m+k}^{(q)}(r)}{P_m'(r)}=Q_{k,q}(r),\qquad
 Q_{0,0}=0,\qquad
 \deg Q_{k,q}\le k+q-1\quad(k+q\ge1).
 \tag{HD8}\label{hd:Qdegree}
\]
Here zero polynomials satisfy every degree upper bound.
They can be constructed without dividing by a polynomial in $r$:
\begin{align}
 Q_{0,1}&=1,&
 Q_{0,q+2}&=\tfrac12\bigl(zQ_{0,q+1}+(q-m)Q_{0,q}\bigr),\notag\\
 Q_{1,q}&=zQ_{0,q}+qQ_{0,q-1}-2Q_{0,q+1},&
 Q_{k+1,q}&=zQ_{k,q}+qQ_{k,q-1}-2(m+k)Q_{k-1,q}\quad(k\ge1).
 \tag{HD9}\label{hd:Qrecursion}
\end{align}
The terms with a negative second index are declared zero.
The first line comes from repeated differentiation of
$2P_m''-zP_m'+mP_m=0$.
The first formula of the second line uses
$P_{m+1}=zP_m-2P_m'$.
The last formula uses the Hermite recurrence and Leibniz' rule.
They prove the asserted root evaluations and the degree bound
by induction. They also show
$Q_{k,q}(-z)=(-1)^{k+q-1}Q_{k,q}(z)$ whenever $k+q\ge1$.

There are consequently polynomials $p_n(z)$, determined recursively
by the original $A_1,\ldots,A_n$, with
\[
 z_j(u)=r_j+\sum_{n\ge1}p_n(r_j)u^n,\qquad
 \deg p_n\le n-1,\qquad
 p_n(-z)=(-1)^{n-1}p_n(z).
 \tag{HD10}\label{hd:pdegree}
\]
These are coefficient identities in the actual convergent implicit
germs; convergence of $p_n(z)$ for unrestricted $z$ is not required.
For an explicit recursion, let $\delta_{n-1}(u,z)=
\sum_{\nu=1}^{n-1}p_\nu(z)u^\nu$. Then
\[
 p_n(z)=-
 \sum_{\substack{0\le k\le n,\ 0\le q\le n-k\\(k,q)\ne(0,1)}}
 \frac{A_kQ_{k,q}(z)}{q!}
      [u^{\,n-k}]\delta_{n-1}(u,z)^q .
 \tag{HD11}\label{hd:precursion}
\]
The convention $\delta^0=1$ includes the term $k=n,q=0$.
To prove it, substitute the previous coefficients into HD6,
Taylor-expand each polynomial at $r_j$, and divide by $P_m'(r_j)$.
The omitted term contributes exactly the unknown $p_n$.
All remaining terms involve only previous coefficients.
A nonzero monomial in $\delta^q$ has indices
$\nu_1+\cdots+\nu_q=n-k$ and degree at most $n-k-q$.
Together with HD8 its total degree is at most $n-1$.
The same count gives parity $n-1$. For $q=0$ the only possible
term is $k=n$, and HD8 gives the same degree and parity.
This proves the induction and every assertion of HD10.

The bound is the crucial extra constraint: the coefficient of
$u^{n+1}$ in an actual root $\theta_j$ is a polynomial of degree
at most $n-1$ in its original leading Hermite root.
It is two degrees below the unrestricted power $r_j^{n+1}$.
Every $A_n$ is allowed its actual value; none has been set to zero.

\HD\textbf{A convergent coordinate built from all actual branches.}
Let $V_r=(r_j^\ell)_{1\le j\le m,\ 0\le\ell<m}$.
Its determinant is the nonzero real Vandermonde product.
Define analytic germs $\beta_\ell(u)$ by the finite, exact system
\[
 \begin{pmatrix}\beta_0(u)\\ \vdots\\ \beta_{m-1}(u)\end{pmatrix}
       =V_r^{-1}
          \begin{pmatrix}\theta_1(u)-ur_1\\
                          \vdots\\\theta_m(u)-ur_m\end{pmatrix}.
 \tag{HD12}\label{hd:finiteinterp}
\]
For each $\ell$ the quotient $\beta_\ell(u)/u^\ell$ is analytic
and even, and has value zero at $u=0$. More precisely,
\[
 \beta_\ell(u)=u^{\ell+2}B_\ell(u^2),\qquad B_\ell\in\mathbb C\{t\}.
 \tag{HD13}\label{hd:divisibility}
\]
For proof, the coefficient of $u^{n+1}$ in HD12 is the coefficient
vector of the degree-less-than-$m$ interpolation polynomial of
$p_n$ at the roots. If $n+1\le\ell$, its degree is at most
$n-1<\ell$, so its $\ell$th coefficient is zero.
There is no $u^0$ or $u^1$ coefficient either.
Thus $\beta_\ell$ vanishes to order at least $\ell+1$.
HD7 implies $\beta_\ell(-u)=(-1)^\ell\beta_\ell(u)$:
substitution $z\mapsto-z$ gives the unique interpolation polynomial
for the permuted data. This improves the order to $\ell+2$.
The remaining even convergent power series is a convergent germ
in $t=u^2$, proving HD13.

We obtain the exact real-coefficient analytic map
\[
 F(t,y)=y+t\sum_{\ell=0}^{m-1}B_\ell(t)y^\ell,\qquad
 F(0,y)=y,\qquad
 F(u^2,ur_j)=\theta_j(u).
 \tag{HD14}\label{hd:F}
\]
For $m\ge2$ its degree in $y$ is less than $m$; for $m=1$ it is
$y+\theta_1(t)$. This convention includes $m=1$ without pretending
that the representative of $y$ in a rank-one quotient is nonzero.
All coefficients of $F$ are convergent, because HD12 is a fixed
finite matrix applied to convergent germs and HD13 removes only
proved powers of $u$. Thus the construction is analytic, not
merely a formal coordinate transformation.

\HD\textbf{The exact algebra isomorphism and the unchanged arithmetic unit.}
Introduce the auxiliary algebra
\[
 A_{\mathrm H}=\mathcal O[y]/\mathcal P_m(t,y),\qquad
 \varphi:A\longrightarrow A_{\mathrm H},\quad \xi\longmapsto F(t,y).
 \tag{HD15}\label{hd:isomorphism}
\]
The map is well-defined: divide the polynomial $W(t,F(t,y))$ by
the monic $\mathcal P_m$. Its remainder has degree less than $m$
and vanishes at all distinct points $y=ur_j$ for $t\ne0$, by HD14.
It is zero there and hence identically zero by analyticity.
Let $T(t)$ be the matrix whose $k$th column is the coefficient
vector of $F(t,y)^k$ modulo $\mathcal P_m$, for $0\le k<m$.
Reduction by a monic polynomial preserves holomorphic coefficients.
At $t=0$, $F(0,y)=y$ and $\mathcal P_m(0,y)=y^m$, so
\[
 T(0)=I_m,\qquad T\in\operatorname{GL}_m(\mathcal O).
 \tag{HD16}\label{hd:T}
\]
It follows that $\varphi$ is an isomorphism of the complete finite
$\mathcal O$ algebras, in their stated original and auxiliary bases.
It is not an assertion that $W$ equals $\mathcal P_m$.

There is also an exact ambient analytic coordinate description.
Since $F_y(0,0)=1$, the analytic inverse-function theorem gives
$y=G(t,\xi)$ with $F(t,G(t,\xi))=\xi$ and $G(t,F(t,y))=y$.
The full factorization reads
\begin{align}
 W(t,F(t,y))&=\mathcal P_m(t,y)V(t,y),& V(0,y)&=1,\notag\\
 H_t(x_0+F(t,y))&=\mathcal U(t,y)\mathcal P_m(t,y),&
 \mathcal U(t,y)&=U(t,F(t,y))V(t,y).
 \tag{HD17}\label{hd:unit}
\end{align}
Here $V$ is the holomorphic quotient and a unit near the origin;
its value at $t=0$ follows by division of $y^m$ by $y^m$.
In particular $\mathcal U(0,y)=\sum_{k\ge0}a_{m+k}y^k$.
The complete original unit, leading coefficient, coordinate,
and all derivatives are present in HD17. The auxiliary polynomial
has acquired an exact isomorphism from the actual algebra; it has
not been substituted for the original arithmetic function.

\HD\textbf{Evaluation of the integral monodromy defect.}
Let $r=\lfloor m/2\rfloor$, $\epsilon=m-2r$, and index positive
Hermite roots by $r_a>0$. Define
\[
 B=\prod_{a=1}^r\mathbb C\{u_a\}\times\mathcal O^\epsilon,\quad
 u_a^2=t,\qquad
 \iota_{\mathrm H}(y)=(r_au_a)_a\times0,\qquad
 \iota=\iota_{\mathrm H}\varphi .
 \tag{HD18}\label{hd:B}
\]
The central component is omitted when $\epsilon=0$.
For $A$, the coordinate values are exactly
$\theta_a(u_a)=F(u_a^2,r_au_a)$ and the central actual branch
$F(t,0)$. Thus $\iota$ is precisely the original branch inclusion.
Its injectivity follows on the punctured disc by distinct evaluations
and then over $\mathcal O$ by freeness.
The sign automorphism $S_B:u_a\mapsto-u_a$ fixes the central factor.
On $A_{\mathrm H}$ it is the regular involution
$S_{\mathrm H}:y\mapsto-y$, since
$\mathcal P_m(t,-y)=(-1)^m\mathcal P_m(t,y)$.
In the auxiliary power basis its matrix is
$P_0=\operatorname{diag}_{0\le k<m}((-1)^k)$.
Consequently
\[
 \boxed{\quad M_A=\varphi^{-1}S_{\mathrm H}\varphi
             =T(t)^{-1}P_0T(t)\in\operatorname{GL}_m(\mathcal O),
 \qquad M_AA=A,\qquad\mathcal D_M=0.\quad}
 \tag{HD19}\label{hd:vanishing}
\]
This agrees with the original parameter holonomy:
both maps exchange the evaluations at $\theta_j(u)$ and
$\theta_{\sigma j}(u)$, so they agree on the etale algebra;
injectivity gives the stated extension on $A$.
Its original special-fibre action is exactly
\[
 M_A(0)(\xi^k)=(-1)^k\xi^k,\qquad
 A/tA=\mathbb C[\xi]/(\xi^m).
 \tag{HD20}\label{hd:central}
\]
Every nilpotent layer is retained. For $m\ge2$ the involution is
nonidentity even though the defect vanishes; for $m=1$ it is
the identity.

In ambient coordinates this involution has the convergent formula
\[
 \Psi(t,\xi)=F(t,-G(t,\xi)),\qquad
 \Psi(t,\Psi(t,\xi))=\xi,\qquad \Psi(0,\xi)=-\xi .
 \tag{HD21}\label{hd:ambient}
\]
It preserves the actual zero divisor with its exact unit:
\[
 H_t(x_0+\Psi(t,\xi))
   =(-1)^m
     \frac{\mathcal U(t,-G(t,\xi))}{\mathcal U(t,G(t,\xi))}
       H_t(x_0+\xi).
 \tag{HD22}\label{hd:fullunit}
\]
This follows by applying HD17 twice and using the parity of
$\mathcal P_m$. The quotient is an analytic unit.
Thus this is a local symmetry of the full original divisor,
with a specified transformation of the original function;
it does not assert that $H_t$ itself has an unweighted reflection
symmetry about the colliding point.

\HD\textbf{The divisibility criterion in the original matrix $J=CD$.}
The previous result also evaluates, entry by entry, the exact
criterion asked for in the original lattice presentation.
For any such branch inclusion write
\[
 J=C(t)D(t),\qquad
 D=\operatorname{diag}(t^{d_k}),\quad d_k=\lfloor k/2\rfloor,
 \qquad K(t)=C(t)^{-1}S_BC(t).
 \tag{HD23}
\]
The matrix $C$ is the entire unit matrix of HC7, not its value at zero.
In its coordinates the defect is
\[
 \mathcal D_M\simeq
 \operatorname{im}\left(
  \mathcal O^m\xrightarrow{\ KD\ }
       \bigoplus_{i=0}^{m-1}\mathcal O/(t^{d_i})\right).
 \tag{HD24}\label{hd:criterion-module}
\]
Indeed $C^{-1}$ sends $J\mathcal O^m$ to $D\mathcal O^m$ and
$S_BJ\mathcal O^m$ to $KD\mathcal O^m$.
It follows that
\[
 M_AA=A
 \ \Longleftrightarrow\
 K_{ik}(t)\in t^{\max(d_i-d_k,0)}\mathcal O
          \quad\hbox{for all }i,k.
 \tag{HD25}\label{hd:criterion}
\]
For necessity and sufficiency, the entries of
$D^{-1}KD$ are $t^{d_k-d_i}K_{ik}$.
Their integrality is exactly the stated criterion; applying the
involution twice makes inclusion equivalent to equality.
Equivalently, every coefficient
$[t^\nu]K_{ik}$ with $0\le\nu<d_i-d_k$ must vanish.

The actual heat equation proves all these conditions simultaneously.
Let $J_{\mathrm H}$ be the inclusion matrix for HD18. Exact evaluation
gives
\[
 J=J_{\mathrm H}T,\qquad
 S_BJ_{\mathrm H}=J_{\mathrm H}P_0,\qquad
 D^{-1}KD=T^{-1}P_0T .
 \tag{HD26}\label{hd:Jidentity}
\]
The right side is integral by HD16; hence every required coefficient
vanishes by a proved identity involving all actual branches.
This evaluates the original defect rather than discarding the
larger quotient $B/A$. That latter quotient, its conductor and
its nilpotent special-fibre map remain those of HC8--13; its length
$\lfloor(m-1)^2/4\rfloor$ can be positive while HD24 is zero.

\HD\textbf{The first arithmetic coefficients and the exact first variation.}
Use the original ratios
\[
 c=A_1,\quad d=A_2,\quad e=A_3,\quad
 b=2d-c^2,\qquad h=e-cd+\frac{c^3}{3}.
 \tag{HD27}
\]
The explicit recursion gives
\[
 p_1=2c,\qquad p_2=bz,\qquad
 p_3=2h\bigl(z^2-2(m+2)\bigr).
 \tag{HD28}\label{hd:firstcoeffs}
\]
For a direct check of the first new even term,
$Q_{0,2}=z/2$,
$Q_{0,3}=z^2/4-(m-1)/2$,
$Q_{1,1}=0$, $Q_{1,2}=m+1$,
$Q_{2,1}=-2(m+2)$ and
$Q_{3,0}=-2z^2+4(m+2)$.
Substitution in HD11 yields respectively the coefficient
$2h$ of $z^2$ and the constant $-4(m+2)h$.
Thus the full branches have
\[
 \theta_j(u)=r_ju+2cu^2+b r_ju^3
             +2h\bigl(r_j^2-2(m+2)\bigr)u^4+O(u^5).
 \tag{HD29}
\]
Every remainder here is from the fixed actual convergent germ.

There is an exact expression for the whole coefficient of $t$,
not just the displayed low powers. Define the original logarithmic
unit coefficients by
\[
 \log\frac{U(0,z)}{a_m}
       =\log\left(1+\sum_{n\ge1}A_nz^n\right)
       =\sum_{n\ge1}\ell_nz^n .
 \tag{HD30}\label{hd:logunit}
\]
This is the unique local logarithm with value zero at zero.
In particular $\ell_1=c$, $\ell_2=d-c^2/2$, $\ell_3=h$, and
\[
 \ell_5=A_5-cA_4-d e+c^2e+cd^2-c^3d+c^5/5.
 \tag{HD31}
\]
For $1\le n\le m$, the coefficient of $z^{n-1}$ in $p_n(z)$
is $2\ell_n$.
Here is a complete coefficient proof. The leading term of
$Q_{0,q}$ for $q\ge1$ is $2^{1-q}z^{q-1}$ by HD9.
The leading term of $Q_{k,0}$ for $k\ge1$ is $-2z^{k-1}$.
For $k,q\ge1$, derivative identities give a representative of degree
at most $k+q-3$: if $q<k$, use
$P_{m+k}^{(q)}=(m+k)!P_{m+k-q}/(m+k-q)!$;
if $q=k$, the value is zero; if $q>k$, use
$P_{m+k}^{(q)}=((m+k)!/m!)P_m^{(q-k)}$.
When $k+q\le n\le m$, this representative and HD9's representative
both have degree less than $m$ and agree at all $m$ roots;
they are therefore the same polynomial. These mixed terms cannot
contribute to degree $n-1$ in HD11 for the stated range.
If $v_n=[z^{n-1}]p_n$ and $V(w)=\sum_{n=1}^m v_nw^n$,
the remaining highest-degree coefficients say exactly
\[
 2\bigl(e^{V(w)/2}-1\bigr)-2\sum_{n=1}^m A_nw^n
          =0\pmod{w^{m+1}}.
 \tag{HD32}
\]
This is an identity of truncated formal coefficient series, so its
unique solution is
$V(w)=2\log(1+\sum A_nw^n)\pmod{w^{m+1}}$.
This proves exactly the coefficients used below, with no
asymptotic argument in $z$ or choice of a polynomial representative
above degree $m-1$.

In HD12--14, a term $t y^{n-1}$ has weight $n+1$ for
$\operatorname{wt}t=2$, $\operatorname{wt}y=1$.
Only the top coefficient of $p_n$ supplies it, and $n-1<m$
means that no interpolation reduction is needed. Hence
\[
 \partial_t F(0,y)=2\sum_{n=1}^m\ell_n y^{n-1},\qquad
 \Psi(t,\xi)=-\xi+
       4t\sum_{\substack{1\le n\le m\\ n\ {\rm odd}}}
                         \ell_n\xi^{n-1}+O(t^2).
 \tag{HD33}\label{hd:firstvariation}
\]
The first equality is an equality of polynomials for the chosen
$F$; it also applies to $m=1$.
The second is an equality of analytic germs in $(t,\xi)$:
differentiate $F(t,-G(t,\xi))$, using
$\partial_tG(0,\xi)=-\partial_tF(0,\xi)$.
Thus the actual unit determines every first-order matrix entry.

For $m\ge2$, in the original basis $1,\xi,\ldots,\xi^{m-1}$,
the first variation is especially explicit. For $1\le k<m$,
\[
 M_A(\xi^k)=(-1)^k\xi^k+
 4k(-1)^{k-1}t
   \sum_{\substack{n\ {\rm odd},\ n\ge1\\k+n-2<m}}
                 \ell_n\,\xi^{k+n-2}+O(t^2A),
 \qquad M_A(1)=1 .
 \tag{HD34}\label{hd:matrixvariation}
\]
Terms already multiplied by $t$ may be reduced modulo $\xi^m$
to this order, since $W(0,\xi)=\xi^m$.
This proves each coefficient and sign in the stated original basis.

\HD\textbf{Multiplicities one and two.}
For $m=1$ there is one analytic branch and $A=B=\mathcal O$.
The action is the identity and the defect is zero.
For $m=2$ there is one quadratic component,
\[
 W=(\xi-E(t))^2-tR(t)^2,\qquad
 M_A(\xi)=2E(t)-\xi,\qquad A=B,\quad\mathcal D_M=0 .
 \tag{HD35}
\]
Here $R(0)=\sqrt2$ and both complete germs are the actual
even and odd branch parts. HD29 evaluates
$E(t)=2ct-12ht^2+O(t^3)$.
The coefficient $-12h$ uses $r_j^2=2$ and $m+2=4$.
Thus $M_A(\xi)=-\xi+4ct-24ht^2+O(t^3)$ exactly in this rank-two
algebra. No unit has been set equal to a constant.

\HD\textbf{Multiplicity three: a nonzero normalization quotient, zero loop defect.}
Write the actual factors as
$Q=(\xi-E(t))^2-tR(t)^2$ and $\xi-\eta(t)$, with $R(0)=\sqrt6$.
All three germs include every higher arithmetic coefficient.
The exact generator image is
\[
 M_A(\xi)=-\xi+2E(t)+2\lambda(t)Q(t,\xi),\qquad
 \lambda(t)=\frac{\eta(t)-E(t)}
                {(\eta(t)-E(t))^2-tR(t)^2}.
 \tag{HD36}\label{hd:m3}
\]
At the two quadratic roots this is $2E-\xi$, and at $\xi=\eta$
it equals $\eta$; injectivity of the full branch evaluations
therefore proves the formula. HD29 gives
\[
 E=2ct-8ht^2+O(t^3),\quad
 \eta=2ct-20ht^2+O(t^3),\quad
 R^2=6(1+2bt)+O(t^2),\quad
 \lambda=2ht+O(t^2).
 \tag{HD37}
\]
In particular the displayed quotient in HD36 is analytic: its
denominator is $t$ times a unit with constant $-6$, and its
numerator has order at least two. This checks directly the first
case with $B/A\simeq\mathcal O/(t)$.
Its entire holonomy matrix begins
\[
 M_A=
 \begin{pmatrix}
  1&4ct&0\\
  0&-1&-8ct\\
  0&4ht&1
 \end{pmatrix}
       +O(t^2),\qquad \mathcal D_M=0 .
 \tag{HD38}
\]
The notation $O(t^2)$ is entrywise and the first column is exactly
$(1,0,0)^T$. The $h$ term is the original coefficient
$a_{m+3}/a_m-cd+c^3/3$, not a freely selected polynomial parameter.

\HD\textbf{Multiplicity four: all two-pair denominators cancel.}
There are two actual quadratic factors
$Q_a=(\xi-E_a)^2-tR_a^2$, with
$r_1^2=6-2\sqrt6$, $r_2^2=6+2\sqrt6$ and $R_a(0)=r_a>0$.
Put, retaining all germs,
\[
 \Delta=E_2-E_1,\quad a=2\Delta,\quad
 b_0=\Delta^2+t(R_2^2-R_1^2),\quad
 \mathfrak r=b_0^2-a^2tR_2^2 .
 \tag{HD39}
\]
Then the exact original-coordinate formula is
\[
 M_A(\xi)=-\xi+2E_1+
     2\Delta Q_1(t,\xi)
       \frac{b_0-a(\xi-E_2)}{\mathfrak r}.
 \tag{HD40}\label{hd:m4}
\]
Indeed modulo $Q_2$,
$Q_1=a(\xi-E_2)+b_0$ and $(\xi-E_2)^2=tR_2^2$.
The fraction is the multiplicative inverse of $Q_1$ modulo $Q_2$.
The added term vanishes modulo $Q_1$ and equals $2\Delta$
modulo $Q_2$, which proves every evaluation in HD40.

HD29 shows $\Delta=2h(r_2^2-r_1^2)t^2+O(t^3)$ and
$b_0=(r_2^2-r_1^2)t+O(t^2)$.
Thus $\mathfrak r=t^2$ times a unit with constant
$(r_2^2-r_1^2)^2=96$. The two coefficient quotients
$\Delta b_0/\mathfrak r$ and $\Delta a/\mathfrak r$ have orders
at least one and two, respectively. This proves analyticity
without cancelling any unspecified arithmetic factor.
Here $B/A\simeq(\mathcal O/(t))^2$, but its loop submodule is zero.
The full matrix has first-order part
\[
 M_A=
 \begin{pmatrix}
 1&4ct&0&0\\
 0&-1&-8ct&0\\
 0&4ht&1&12ct\\
 0&0&-8ht&-1
 \end{pmatrix}
       +O(t^2),\qquad \mathcal D_M=0 .
 \tag{HD41}
\]
Again every higher coefficient is specified by the convergent
original-coordinate formula HD40 or, equivalently, HD19.

\HD\textbf{The next multiplicities and the first additional constraints.}
For $m=5,6$ the largest exponent $d_i$ in HD23 is two, so
$t^2B\subset A$. All actual pair centres and the central branch
have common term $2ct$ by HD29. The even-centre vector is therefore
$2ct\cdot1+t^2b_2$ with $b_2\in B$, and lies in $A$.
The generator image $M_A(\xi)$ is twice this vector minus $\xi$,
which directly proves the criterion in these two cases.
The first new coefficient of its matrix is completely determined:
\[
 M_A(\xi)=-\xi+4t(c+h\xi^2+\ell_5\xi^4)+O(t^2A),
       \qquad m=5\ \hbox{or}\ 6 .
 \tag{HD42}
\]
All entries of the remaining columns are given by HD34.
The coefficient $\ell_5$ is the full arithmetic expression HD31.

For $m=7,8$ one has $t^3B\subset A$. Now the exact fourth-order
root coefficient is needed. The pair-centre vector, including
the central component if present, satisfies
\[
 E_{\mathrm{vec}}=
  2ct\cdot1+
  2ht(\xi-2ct)^2-4(m+2)ht^2\cdot1\pmod{t^3B}.
 \tag{HD43}
\]
To check both kinds of $B$ coordinates, on each pair
$\xi-2ct=r_au_a+b r_au_a^3+O(u_a^4)$.
Its even square is $r_a^2t+O(t^2)$.
Its odd square has order at least $u_a^5$; after division by
$u_a$ and multiplication by $t$ it belongs to $t^3\mathcal O$.
Thus the right side has the required even coefficient
$2ct+2h(r_a^2-2(m+2))t^2$, with zero odd coefficient modulo
$t^3B$. The central coordinate gives
$2ct-4(m+2)ht^2$ as required. This proves HD43.
Every displayed term lies in $A$, so the centre vector does too,
and again $\mathcal D_M=0$.
For all larger multiplicities HD10--19 give the finite, convergent
continuation of these cancellations, rather than a finite-order
presumption about uncomputed terms.

\HD\textbf{Retained fibres and the exact additional receiver.}
The vanishing of HD4 concerns this specific loop-lattice defect.
It does not annihilate $B/A$, its conductor, the cotangent torsion,
or the original trace radical. For example HD20 acts by signs
on all powers in $\mathbb C[\xi]/\xi^m$, while the original
reflected trace at the endpoint remains
$m\overline{f(0)}g(0)$ and has radical $(\xi)$ when $m>1$.
The endpoint $v=i\xi/2$ is the original arithmetic jet coordinate,
and the loop acts by $v\mapsto-v$ there. The arithmetic multiplier
$e^{iLx_0/2}\exp(LM_v)$ remains its separate full operator;
the commutator and trace comparisons proved in HH33--38 apply
to the now established regular lattice automorphism.

For every original bounded distributive support lattice $\Lambda$,
there is now an actual lifted algebra involution
\[
 G_\Lambda(A)\longrightarrow G_\Lambda(A),\qquad
 (a,1_\Lambda)\longmapsto(M_Aa,1_\Lambda),\qquad
 (0,\lambda)\longmapsto(0,\lambda).
 \tag{HD44}
\]
It preserves both operations because $M_A$ is a unital
$\mathcal O$-algebra automorphism and every support coordinate
is fixed. Its square is the identity; its arithmetic projection
is HD19 and its support projection is the identity on the
complete original lattice. This proves the strongest extension
specified by the formerly unevaluated defect.
It supplies no nonreal zero of the original $H_0$ and no negative
test for the complete original Weil form.

